\documentclass{overview}
\title{Calculus Overview - Series}

\newcommand{\sumOInf}{\sum_{n=0}^{\infty}}

\begin{document}
\maketitle

\section{Convergence}

\begin{lemma}[Abel Transformation]
	Let $\{a_n\}$ and $\{b_n\}$ be two sequences, and $A_n = \sum_{k=0}^{n} a_k$, then
	\[
		\sum_{n=0}^{N} a_n b_n = A_N b_N - \sum_{n=0}^{N-1} A_n (b_{n+1} - b_n),
	\]
	\begin{remark}
		A simple way to remember this is to think of the summation as a discrete version
		of integration by parts.
		\[
		\sum b_n \Delta A_n = A_n b_n - \sum A_n \Delta b_n
		\] 
	\end{remark}
\end{lemma}
\begin{theorem}[Abel's Theorem for power series]
	If the power series $\sum_{n=0}^{\infty} a_n x^n$ converges at $x = x_0$, then
	it absolutely converges for all $x$ such that $|x| < |x_0|$.
	and if it diverges at $x = x_0$, then it diverges for all $x$ such that $|x| > |x_0|$.
\end{theorem}
\begin{proof}
	We first prove the first part. Let $b_n=a_nx_0^n$ then $\sum_{n=0}^{\infty} b_n$ converges,
		implying that $\{b_n\}$ is bounded, i.e. there exists $M>0$ such that $|b_n|\leq M$ for all $n$.
	Let $x=tx_0$, then $|t|<1$ and we have
	\[
		\sum_{n=0}^{\infty} |a_n x^n| = \sum_{n=0}^{\infty} |b_n| |t|^n 
		\leq M \sum_{n=0}^{\infty} |t|^n
	\]
	which is a convergent geometric series since $|t|<1$. 
	Hence $\sum_{n=0}^{\infty} a_n x^n$ converges absolutely for $|x|<|x_0|$.
	
	The second part can be easily proved by contraposition.
\end{proof}

\begin{theorem}[Inference]
	If a power series $\sum_{n=0}^{\infty} a_n x^n$
	\begin{itemize}
	\item converges at some point $x_0$, but
	diverges at $-x_0$ then the radius of convergence is exactly $|x_0|$.
	\item conditionally converges at some point
	$x_0$, then the radius of convergence is exactly $|x_0|$.
	\end{itemize}
	\begin{remark}
		For general power series $\sum_{n=0}^{\infty} a_n (x-x_0)^n$, the conclusion
		is become the end points of the convergence interval instead of directly the radius.
	\end{remark}
\end{theorem}

\begin{example}
	If the series $\sum_{n=0}^{\infty} a_n (x+2)^n$ converges at $x=0$ and diverges at
	 $x=-4$, then the radius of convergence is $2$.
\end{example}

\begin{example}
	If the series $\sum_{n=0}^{\infty} a_n$ is conditionally convergent, try to 
	determine $x=\sqrt{3}$ and $x=-\sqrt{3}$ are in the interval of convergence for
	$\sum_{n=0}^{\infty} a_n (x-1)^n$ or not.
	\begin{answer}
		The conditional convergence tells us that the series $\sum_{n=0}^{\infty} a_nx^n$
		has radius of convergence $1$, so the series $\sum_{n=0}^{\infty} a_n (x-1)^n$ has
		convergence interval $[0,2]$. Hence $x=\sqrt{3}$ is in the interval of convergence, but
		$x=-\sqrt{3}$ is not.
	\end{answer}
\end{example}


\begin{theorem}[d'Alembert's Ratio Test]
	Given a series $\sum_{n=0}^{\infty} a_n$, if the limit
	\[
		L=\lim_{n\to\infty} \left| \frac{a_{n+1}}{a_n} \right|
	\] 
	exists, then
	\begin{itemize}
		\item If $L<1$, then the series converges absolutely.
		\item If $L>1$, then the series diverges.
		\item If $L=1$, then the test is inconclusive.
	\end{itemize}
\end{theorem}
\begin{remark}
	Usually used for series with factorials $n!$ or exponentials $a^n$.
\end{remark}
\begin{theorem}[Cauchy's Root Test]
	Given a series $\sum_{n=0}^{\infty} a_n$, if the limit
	\[
		L=\lim_{n\to\infty} \sqrt[n]{|a_n|}
	\] 
	exists, then
	\begin{itemize}
		\item If $L<1$, then the series converges absolutely.
		\item If $L>1$, then the series diverges.
		\item If $L=1$, then the test is inconclusive.
	\end{itemize}
\end{theorem}
\begin{remark}
	Usually used for complex power with $n$ in the exponent like $(n / n+1)^n$.
\end{remark}

\subsection{Properties of Convergent Sets}
The convergent set of linear combinations of convergent series is the intersection
of the convergent sets of the individual series.

\begin{theorem}[Cauchy-Hadamard]
	\label{thm:cauchy-hadamard}
	The radius of convergence $R$ of the power series $\sum_{n=0}^{\infty} a_n x^n$ is given by
	\[
		\frac{1}{R} = \limsup_{n\to\infty} \sqrt[n]{|a_n|}
	\] 
\end{theorem}
\begin{theorem}
	Given a power series $\sum_{n=0}^{\infty} a_n x^n$ with radius of convergence $R$, then
	muplying the series coefficient by any polynomial $P(n)$ does not change the radius of convergence.
\end{theorem}
\begin{proof}
	In fact multiplying the series coefficient by
	any good function $f(n) = \exp(o(n))$ will not change the radius of convergence,
	since
	\[
		\limsup_{n\to\infty} \sqrt[n]{|f(n)|} = \limsup_{n\to\infty} \exp(o(n)/n) = 1
	\] 
	by \ref{thm:cauchy-hadamard}, we have new radius of convergence $R'$ given by
	\[
		\frac{1}{R'} = \limsup_{n\to\infty} \sqrt[n]{|P(n) a_n|} 
		= \limsup_{n\to\infty} \sqrt[n]{|P(n)|} \cdot \limsup_{n\to\infty} \sqrt[n]{|a_n|}
		= \frac{1}{R}
	\] 
\end{proof}
\begin{remark}
	That means when determining the radius of convergence for power series, we can
	ignore any $\exp(o(n))$ factors.
\end{remark}
\begin{example}
	Knowning that the series $\sumOInf a_n x^n$ has radius of convergence $3$ then
	the series $\sumOInf n a_n (x-1)^{n+1}$ has radius of convergence $3$ as well.
\end{example}

\begin{theorem}[Subsequence Convergence]
	If a series $\sumOInf a_n x^n$ has radius of convergence $R$, then any subsequence
	power series $\sumOInf a_{n_k} x^{n_k}$ has radius of convergence at least $R$.
\begin{proof}
	For any subsequence $\{n_k\}$, we have (by the definition of $\limsup$)
	\[
		\limsup_{k\to\infty} |a_{n_k}|^{1/n_k} \leq \limsup_{n\to\infty} |a_n|^{1/n}
	\] 
	then \ref{thm:cauchy-hadamard} finishes the proof.
\end{proof}
\end{theorem}
\begin{remark}
	Be careful that this theorem requires the exponents of $x$ has the same
	subsequence as the coefficients.
\end{remark}

\section{Taylor Series}

\[
	e^x = \sum_{n=0}^{\infty} \frac{x^n}{n!} 
	= 1 + x + \frac{x^2}{2!} + \frac{x^3}{3!} + \cdots
\] 
\[
	(1+x)^{\alpha} = \sum_{n=0}^{\infty} \binom{\alpha}{n} x^n
	= \sum_{n=0}^{\infty} \frac{\alpha(\alpha-1)\cdots(\alpha-n+1)}{n!} x^n
	= 1 + \alpha x + \frac{\alpha(\alpha-1)}{2!} x^2 + \cdots
\] 
\[
	\sin x = \sum_{n=0}^{\infty} \frac{(-1)^n}{(2n+1)!} x^{2n+1}
	= x - \frac{x^3}{3!} + \frac{x^5}{5!} - \frac{x^7}{7!} + \cdots
\] 
\[
	\cos x = \sum_{n=0}^{\infty} \frac{(-1)^n}{(2n)!} x^{2n}
	= 1 - \frac{x^2}{2!} + \frac{x^4}{4!} - \frac{x^6}{6!} + \cdots
\] 
\subsection{Convert Basic Series to Other Series}
The basic series geometric series
\begin{equation}\label{eq:geometric-series}
	\frac{1}{1-x} = \sum_{n=0}^{\infty} x^n
\end{equation}
take the integral of both sides gives
\begin{equation} \label{eq:log-series-neg}
	-\ln(1-x) = \sum_{n=0}^{\infty} \frac{x^{n+1}}{n+1} =
 \sum_{n=1}^{\infty} \frac{x^n}{n}
\end{equation}
\begin{equation} \label{eq:log-series}
	\ln(1+x) = \sum_{n=1}^{\infty} \frac{(-1)^{\textcolor{red}{n-1}}}{n} x^n
	= x - \frac{x^2}{2} + \frac{x^3}{3} - \frac{x^4}{4} + \cdots
	\quad x \in (-1,1]
\end{equation}
calculating \eqref{eq:log-series} add \eqref{eq:log-series-neg} gives
\begin{equation}
	\sum_{n=1}^{\infty} \frac{x^{2n-1}}{2n-1} = \frac{1}{2}\left(
		\sum_{n=1}^{\infty} \frac{x^n}{n} + \sum_{n=1}^{\infty} \frac{(-x)^{n-1}}{n}
	\right) 
	= \frac{-\ln(1-x) + \ln(1+x)}{2}
	= \frac{1}{2} \ln\left(\frac{1+x}{1-x}\right)
\end{equation}
extract the even terms of \eqref{eq:log-series-neg} gives
\begin{equation}
	-\frac{1}{2} \ln(1-x^2) =
	\frac{1}{2} \sum_{n=1}^{\infty} \frac{(x^2)^n}{n}
	= \sum_{n=1}^{\infty} \frac{x^{2n}}{2n}
\end{equation}
substituting $x$ with $-x^2$ in \eqref{eq:geometric-series} gives
\begin{equation}
	\frac{1}{1+x^2} = \sum_{n=0}^{\infty} (-1)^n x^{2n}
\end{equation}
take the integral of both sides gives
\begin{equation}
	\atan x = \sum_{n=0}^{\infty} (-1)^n \frac{x^{2n+1}}{2n+1}
	= \sum_{n=1}^{\infty}  (-1)^{n-1}\frac{x^{2n-1}}{2n-1}
\end{equation}
take the derivative of both sides of \eqref{eq:geometric-series} gives
\begin{equation}
	\frac{1}{(1-x)^2} = \sum_{n=0}^{\infty} (n+1) x^n = \sum_{n=1}^{\infty} n x^{n-1}
\end{equation}
take the derivative of both sides again gives
\begin{equation}
	\frac{2}{(1-x)^3} = \sum_{n=0}^{\infty} (n+1)(n+2) x^n = \sum_{n=2}^{\infty} n(n-1) x^{n-2}
\end{equation}

\subsection{Manipulate and complete series terms to get known series}

In order to get a known series list above, we notice that most of the known series
require the series terms has the form of
\[
	\frac{x^\square}{\square} \quad
	\text{or} \quad
	\frac{x^\square}{\square!}
\] 
which we call it the denominator form. Or another form
\[
	n^{\underline{k}} x^{n-k}
\] 
where $n^{\underline{k}} = n(n-1)(n-2)\cdots(n-k+1)$ is the falling factorial, 
which we call it the numerator form.

We want to manipulate the series terms to get close to the known series form.
For example, like rearranging
\[
	\begin{aligned}
		\sum (an^2 + bn + c) x^n 
		&= \sum (a n(n-1) + (a+b)n + c) x^n\\
		&= \frac{2ax^2}{(1-x)^3} + \frac{(a+b)x}{(1-x)^2} + \frac{c}{1-x}
	\end{aligned}
\] 
\begin{example}
	Find the sum function of the series
	\[
		\sum_{n=1}^{\infty} \frac{1}{n 2^n} x^{n+1}
	\] 
	\begin{answer}
		\[
			S(x) = x \sum_{n=1}^{\infty} \frac{1}{n} \left(\frac{x}{2}\right)^n
			= -x \ln\left(1-\frac{x}{2}\right)
		\] 
	\end{answer}
\end{example}
\begin{example}
	Find the value of the series
	\[
		\sum_{n=2}^{\infty} \frac{1}{(n^2-1) 2^n}
	\]
	\begin{answer}
		\[
			\begin{aligned}
				S(x) &= \sum_{n=2}^{\infty} \frac{1}{n^2-1} x^n 
				= \frac{1}{2} \sum_{n=2}^{\infty} \frac{1}{n-1} x^n
				- \frac{1}{2} \sum_{n=2}^{\infty} \frac{1}{n+1} x^n\\
						 &= \frac{x}{2} \sum_{n=1}^{\infty} \frac{x^{n}}{n}
						 - \frac{1}{2x} \sum_{n=3}^{\infty} \frac{x^{n}}{n}\\
						 &= -\frac{x}{2} \ln(1-x) - \frac{1}{2x} \left(-\ln(1-x) - x - \frac{x^2}{2}\right)\\
						 &= \frac{1-x^2}{2x} \ln(1-x) + \frac{x}{4} + \frac{1}{2}
			\end{aligned}
		\] 
		substituting $x=\frac{1}{2}$ gives
		\[
			S\left(\frac{1}{2}\right) = \frac{5}{8} -\frac{3}{4} \ln 2
		\]
	\end{answer}
\end{example}

\end{document}
